% Gauss, Werke X/1 (1917), native leaves45–46, printed38–39.
\SourceLeaf{045}\NoSourceHead
\SourceGroup{VI}
\HistoricTitle{[ZUR LEHRE VON DEN GANZEN KOMPLEXEN ZAHLEN.]}
\HistoricCitation{[Ein Zettel in Ec2, Kapsel 43.]}
\NumberedItem{1}
\par\indent\textit{Sind $a,b$ Primzahlen unter sich, so werden entweder beide $ax+by$, $ay-bx$ oder keine durch $aa+bb$ theilbar sein.}

\textit{Bew[eis.]}

Multiplicirt man die beiden identischen Gleichungen
\[
\begin{aligned}
(a-bi)(x+iy)&=ax+by+(ay-bx)i\\
(a+bi)(\alpha-\beta i)&=a\alpha+b\beta+(b\alpha-a\beta)i
\end{aligned}
\]
\noindent in einander, so wird, identisch,
\[
\begin{gathered}
(aa+bb)\{\alpha x+\beta y+(\alpha y-\beta x)i\}=(a\alpha+b\beta)(ax+by)+\\
+(a\beta-b\alpha)(ay-bx)+\{(a\alpha+b\beta)(ay-bx)+(ax+by)(b\alpha-a\beta)\}i.
\end{gathered}
\]
\par\indent Bestimmt man also $\alpha,\beta$ so, dass
\[
a\alpha+b\beta=1,
\]
\noindent so wird
\[
\begin{aligned}
(aa+bb)(\alpha x+\beta y)+(b\alpha-a\beta)(ay-bx)&=ax+by\\
(aa+bb)(\alpha y-\beta x)-(b\alpha-a\beta)(ax+by)&=ay-bx.
\end{aligned}
\]
\par\indent Die $\begin{pmatrix}\text{erste}\\\text{zweite}\end{pmatrix}$ Gleichung beweiset, dass, wenn $\begin{pmatrix}ay-bx\\ax+by\end{pmatrix}$ durch $aa+bb$ theilbar ist, auch $\begin{pmatrix}ax+by\\ay-bx\end{pmatrix}$ es sein werde.
\clearpage

\SourceLeaf{046}\ArithmeticHead{39}{R}{ÜBER GANZE KOMPLEXE ZAHLEN.}
\NumberedItem{2}
\par\indent Ist eine Primzahl
\[
p=3n+1=x^2+3y^2,
\]
\noindent so ist
\[
2x\equiv\frac{\tfrac32n(\tfrac32n-1)\ldots(n+1)}{1\ldots\tfrac12n}\qquad(\mathrm{m[od.]}\ p)
\]
\par\indent Unter derselben Voraussetzung ist
\[
4p=l^2+27m^2
\]
\noindent und
\[
l\equiv\frac{-(n+1)(n+2)\ldots2n}{1.2\ldots n}.
\]
\par\indent Ist eine Primzahl
\[
p=8n+1=a^2+2b^2,
\]
\noindent so ist
\[
2a\equiv\frac{4n\ldots(3n+1)}{1\ldots n};
\]
\noindent ist
\[
p=8n+3=a^2+2b^2,
\]
\noindent so ist
\[
2a\equiv\frac{(4n+1)\,4n\ldots(3n+2)}{1\ldots n}.
\]
\ClosingRule
\NumberedItem{3}
\HistoricCitation{[Zettel in Ec3, Kapsel 43]}
\par\indent Der Modulus $a+ib$ und $aa+bb=D$.

Kleinster Rest von $x+iy\ (\mathrm{mod.}\ a+bi)$
\[
x+iy-(a+ib)\left\{\left[\frac{ax+by}{D}\right]+i\left[\frac{ay-bx}{D}\right]\right\}.
\]
\clearpage
